\subsection{\texorpdfstring{Criterios de coherencia para los módulos
$\mathcal H^0_{X/Z}(\mathcal F)$.}{Criterios de coherencia para los módulos H\string^0\string_\string{X/Z\string}(F)}}

\begin{proposition}[5.11.1]
\label{IV.2.s5.11.1-fr}
Sean $X$ un preesquema localmente noetheriano, $Z$ una parte de $X$ estable por
especialización, $\mathcal F$ un $\mathcal O_X$-módulo coherente. Designemos por
$(x_\alpha)$ la familia de los puntos de
$\operatorname{Ass}(\mathcal F)\cap(X\setmin Z)$ y, para todo $\alpha$, sean
$Y_\alpha$ el subpreesquema cerrado reducido de $X$ que tiene
$\overline{\{x_\alpha\}}$ por espacio subyacente,
$Z_\alpha=Z\cap\overline{\{x_\alpha\}}$. Entonces las dos condiciones siguientes
son equivalentes:
\begin{enumerate}
\item[\textup{a)}] $\mathcal H^0_{X/Z}(\mathcal F)$ es un $\mathcal O_X$-módulo
coherente.
\item[\textup{b)}] Para todo $\alpha$,
$\mathcal H^0_{Y_\alpha/Z_\alpha}(\mathcal O_{Y_\alpha})$ es un
$\mathcal O_{Y_\alpha}$-módulo coherente.
\end{enumerate}
Estas dos condiciones implican la siguiente:
\begin{enumerate}
\item[\textup{c)}] Para todo $\alpha$, se tiene
$\operatorname{codim}(Z_\alpha,Y_\alpha)\geq2$.
\end{enumerate}
Además, las tres condiciones \textup{a), b), c)} son equivalentes cuando se verifica
asimismo \emph{una} de las propiedades siguientes:
\begin{enumerate}
\item[\textup{(i)}] Todo punto de $X$ admite una vecindad abierta isomorfa a un
subesquema de un esquema regular \textup{(en cuyo caso se dice también que $X$ es}
localmente inmersible en un esquema regular\textup{)}.
\item[\textup{(ii)}] Para todo $\alpha$, $Y_\alpha$ es universalmente catenario
\textup{(5.6.2)} y su normalización \textup{(\textbf{II}, 6.3.8)} $Y'_\alpha$ es finita sobre $Y_\alpha$.
\end{enumerate}
\end{proposition}

Todas las propiedades consideradas son locales sobre $X$, por lo que es posible
limitarse al caso en que $X=\operatorname{Spec}(A)$ es afín, siendo $A$ un
anillo noetheriano, y $\mathcal F=\widetilde M$, donde $M$ es un $A$-módulo
de tipo finito. Entonces, para todo $\alpha$, si $h_\alpha$ es la inyección
canónica $Y_\alpha\to X$,
$(h_\alpha)_*(\mathcal O_{Y_\alpha})=\mathcal G_\alpha$ es el
$\mathcal O_X$-módulo correspondiente al $A$-módulo cociente
$A/\mathfrak j_{x_\alpha}$, y, por definición de
$\operatorname{Ass}(\mathcal F)$, este $A$-módulo es isomorfo a un
sub-$A$-módulo de $M$. Como $\mathcal H^0_{X/Z}(\mathcal G_\alpha)$ es un
$\mathcal O_X$-submódulo cuasicoherente de
$\mathcal H^0_{X/Z}(\mathcal F)$ (5.9.2), la hipótesis de que
$\mathcal H^0_{X/Z}(\mathcal F)$ sea coherente implica que lo mismo ocurre con
$\mathcal H^0_{X/Z}(\mathcal G_\alpha)$. Por otra parte, de la
definición (5.9.1.2) resulta que $\mathcal H^0_{X/Z}(\mathcal G_\alpha)$ es
isomorfo a
$\mathcal H^0_{Y_\alpha/Z_\alpha}(\mathcal O_{Y_\alpha})$;
esto prueba que \emph{a)} implica \emph{b)}. Para ver que \emph{b)} implica \emph{a)}, basta demostrar que
existe una filtración finita $(\mathcal F_i)_{0\leq i\leq n}$ de
$\mathcal F$ formada por $\mathcal O_X$-módulos coherentes, con
$\mathcal F_0=\mathcal F$, $\mathcal F_n=0$, y tal que
$\mathcal H^0_{X/Z}(\mathcal F_i/\mathcal F_{i+1})$ sea coherente para todo
$i$: esto resulta, por inducción descendente sobre $i$, de las sucesiones
exactas
\[
0\longrightarrow\mathcal F_{i+1}\longrightarrow\mathcal F_i
\longrightarrow\mathcal F_i/\mathcal F_{i+1},
\]
del hecho de que $\mathcal H^0_{X/Z}$ es un funtor exacto por la izquierda
(5.9.2), y finalmente de (\textbf 0$_\mathrm I$, 5.3.3) y (\textbf I, 6.1.1). En virtud de
(3.2.8), basta, por tanto, probar que
$\mathcal H^0_{X/Z}(\mathcal F)$ es coherente cuando $\mathcal F$ es
\emph{irredundante}; dicho de otro modo,
$\operatorname{Ass}(M)=\{\mathfrak p\}$ consta de un solo elemento. Señalemos
ahora el

\begin{lemma}[5.11.1.1]
\label{IV.2.s5.11.1.1-fr}
Sean $A$ un anillo noetheriano, $M$ un $A$-módulo de tipo finito, tal que
$\operatorname{Ass}(M)=\{\mathfrak p\}$. Existe una filtración finita
$(M_h)_{0\leq h\leq m}$ de $M$ tal que $M_0=M$, $M_m=0$ y que
$M_h/M_{h+1}$ sea isomorfo a un submódulo de $A/\mathfrak p$.
\end{lemma}

Señalemos primero que el homomorfismo canónico $M\to M_\mathfrak p=N$ es
inyectivo (Bourbaki, \emph{Alg. comm.}, cap. IV, \S~1, n\textsuperscript{o}~2,
prop. 6). Pongamos $B=A_\mathfrak p$, $\mathfrak m=\mathfrak pA_\mathfrak p$,
ideal maximal de $B$; se tiene $\operatorname{Ass}(N)=\{\mathfrak m\}$
(\emph{loc. cit.}, prop. 5), y como $N$ es un $B$-módulo de tipo finito,

\oldpage[IV-2]{123}

existe un entero $r$ tal que $\mathfrak m^rN=0$; si se pone
$N'_j=\mathfrak m^jN$ para $0\leq j\leq r$, $N'_j/N'_{j+1}$ es un
$(B/\mathfrak m)$-módulo de tipo finito, por tanto suma directa de un número
finito de submódulos isomorfos a $B/\mathfrak m$, puesto que
$B/\mathfrak m$ es un cuerpo; dicho de otro modo, existe una filtración finita
$(N_h)_{0\leq h\leq m}$ de $N$ tal que $N_m=0$ y que $N_h/N_{h+1}$ sea
isomorfo a $B/\mathfrak m$, es decir, al cuerpo de fracciones de
$A/\mathfrak p$; la filtración formada por los $M_h=M\cap N_h$ cumple lo pedido, pues
$M_h/M_{h+1}$ es isomorfo a un sub-$(A/\mathfrak p)$-módulo de tipo finito de
$N_h/N_{h+1}=B/\mathfrak m$; pero se sabe que tal submódulo es isomorfo a un
submódulo de $A/\mathfrak p$.

La existencia de la filtración $(M_h)$ muestra entonces, por el mismo
razonamiento que antes, que para probar (bajo \emph{b)}) que
$\mathcal H^0_{X/Z}(\mathcal F)$ es coherente, es posible limitarse a demostrar
que $\mathcal H^0_{X/Z}(\mathcal P)$ es coherente, donde
$\mathcal P=(A/\mathfrak p)^\sim$, siendo $\mathfrak p$ un ideal asociado a
$M$. Pero si $\mathfrak p=\mathfrak j_y$ con $y\in Z$, el soporte de
$\mathcal P$ es un conjunto cerrado contenido en $Z$, puesto que $Z$ es
estable por especialización; la definición (5.9.1.2) muestra entonces que
$\mathcal H^0_{X/Z}(\mathcal P)=0$. Si, por el contrario,
$\mathfrak p=\mathfrak j_{x_\alpha}$ para un $\alpha$, se tiene
$\mathcal P=(h_\alpha)_*(\mathcal O_{Y_\alpha})$ por definición, y ya se ha
visto que $\mathcal H^0_{X/Z}(\mathcal P)$ es isomorfo a
$\mathcal H^0_{Y_\alpha/Z_\alpha}(\mathcal O_{Y_\alpha})$,
por tanto es coherente en virtud de la hipótesis \emph{b)}.

Ya se ha visto (5.10.10, (ii)) que \emph{a)} implica \emph{c)}. Queda probar que, bajo una
u otra de las hipótesis (i), (ii), \emph{c)} implica \emph{b)}. Observemos que, si $X$
verifica (i), lo mismo ocurre con cada $Y_\alpha$. Basta, pues, en virtud de
(5.9.3.1), probar el

\begin{corollary}[5.11.2]
\EGAReaderTarget{ega4.c04.cor.5-11-2}{IV.2.s5.11.2-fr}
Sea $A$ un anillo noetheriano íntegro que verifica una de las dos hipótesis
siguientes:
\begin{enumerate}
\item[\textup{(i)}] $A$ es cociente de un anillo noetheriano regular.
\item[\textup{(ii)}] $A$ es universalmente catenario, y su clausura integral $A'$ es
una $A$-álgebra finita.
\end{enumerate}
Entonces el anillo $A^{(1)}$, intersección de los $A_\mathfrak p$ donde
$\mathfrak p$ recorre el conjunto de los ideales primos de $A$ de altura 1,
es una $A$-álgebra finita.
\end{corollary}

(i) Pongamos $X=\operatorname{Spec}(A)$. El conjunto $U$ de los puntos
$x\in X$ en los que $X$ verifica $(S_2)$ es \emph{abierto} en $X$ en virtud de la
hipótesis (i) (6.11.2)
\footnote{El lector podrá comprobar que (5.11.2) no se utiliza en la
demostración de (6.11.2).}. Además, si se pone $Z=X\setmin U$, se tiene
$\operatorname{codim}(Z,X)\geq2$: en efecto, para todo $x\in X$ tal que
$\dim(A_x)\leq1$, se tiene $\dim(A_{x'})\leq1$ para todo $x'$ generización de
$x$, y como $A_{x'}$ es íntegro, $\operatorname{prof}(A_{x'})\geq1$, por tanto
$X$ verifica $(S_2)$ en el punto $x$. Se tiene, pues, $Z\subset Z^{(2)}$, y
$\mathcal H^0_{X/Z^{(2)}}(\mathcal O_X)
=j_*(\mathcal H^0_{U/Z^{(2)}}(\mathcal O_U))$, donde $j:U\to X$ es la
inyección canónica (5.9.1.2). Pero como el preesquema $U$ verifica $(S_2)$,
resulta de (5.10.14) que
$\mathcal H^0_{U/Z^{(2)}}(\mathcal O_U)$ es isomorfo a $\mathcal O_U$; por
otra parte, como $\operatorname{codim}(Z,X)\geq2$, se sabe, según el capítulo
III, \S~9, que $j_*(\mathcal O_U)$ es un $\mathcal O_X$-módulo coherente, lo
que demuestra la proposición en el caso (i).

(ii) El anillo $A'$ es, en virtud de la hipótesis (ii), un anillo \emph{noetheriano}
íntegro e integralmente cerrado, por tanto (Bourbaki, \emph{Alg. comm.}, cap.
VII, \S~1, n\textsuperscript{o}~6, th. 4) intersección de sus anillos locales
$A'_{\mathfrak p'}$, donde $\mathfrak p'$ recorre el conjunto de los ideales
primos de altura 1 de $A'$. Ahora bien, para tal ideal primo $\mathfrak p'$,
si se pone $\mathfrak p=\mathfrak p'\cap A$ y $S=A\setmin\mathfrak p$,
$A'_{\mathfrak p'}$ es un anillo local en el ideal primo
$S^{-1}\mathfrak p'$ de $S^{-1}A'$, y $S^{-1}A'$ es por hipótesis

\oldpage[IV-2]{124}

una $A_\mathfrak p$-álgebra finita; como $S^{-1}\mathfrak p'$ está por encima
del ideal maximal $\mathfrak pA_\mathfrak p$ de $A_\mathfrak p$, es un ideal
maximal de $S^{-1}A'$; pero en virtud de la hipótesis (ii) y de (5.6.3, (i)),
$A_\mathfrak p$ es universalmente catenario. Se deduce, pues, de (5.6.10) que
$\dim(A_\mathfrak p)=\dim(A'_{\mathfrak p'})=1$. De ello resulta que se tiene
$A^{(1)}\subset A'$, y como $A'$ es un $A$-módulo de tipo finito por hipótesis,
lo mismo ocurre con $A^{(1)}$, puesto que $A$ es noetheriano; la proposición
queda, por tanto, demostrada en el caso (ii).

\begin{remark}[5.11.3]
\label{IV.2.s5.11.3-fr}
El hecho de que $A^{(1)}$ sea una $A$-álgebra finita ya no es necesariamente
cierto si, en la hipótesis (ii) de (5.11.2), solo se supone que $A$ es
\emph{catenario}. Un ejemplo lo proporciona el anillo local catenario $A$ construido
en (5.6.11), cuya clausura integral $A'$ (denotada $B$ en (5.6.11)) es una
$A$-álgebra finita; si $A^{(1)}$ fuese también una $A$-álgebra finita, como
está contenido en el cuerpo de fracciones de $A$, estaría contenido en $A'$.
Pero, por otra parte, con las notaciones de (5.6.11), todo ideal primo de
altura 1 de $A$ es de la forma $\mathfrak pA$, donde $\mathfrak p$ es un ideal
primo de altura 1 de $C$, y se tiene
$A_{\mathfrak pA}=C_\mathfrak p$; se sabe (5.6.11) que
$C_\mathfrak p=E_{\mathfrak p'}$, donde $\mathfrak p'$ es el único ideal
primo de $E$ por encima de $\mathfrak p$, por tanto $A_{\mathfrak pA}$ es
integralmente cerrado y contiene, en consecuencia, $A'$; por definición, se
tiene, pues, $A'\subset A^{(1)}$, y la hipótesis de que $A^{(1)}$ sea una
$A$-álgebra finita implicaría finalmente $A^{(1)}=A'$. Pero esta conclusión
es absurda, pues uno de los dos ideales primos de $A'$ por encima del ideal
maximal $\mathfrak nA$ de $A$ es de altura 1, mientras que $\mathfrak nA$ es
de altura 2, lo que contradiría (5.10.17, (iv)).
\end{remark}
